%% test-breakable-long.tex
%% =======================
%% Typesetting test for every engine: breakable boxes with "page
%% footnote" that TeX breaks across pages (no \tcbbreak).
%%
%%   pdflatex test-breakable-long.tex   (or xelatex / lualatex)
%%   platex  test-breakable-long.tex && dvipdfmx test-breakable-long
%%   uplatex -jobname=test-breakable-long "\def\CLS{ujarticle}\input test-breakable-long"
%% The class is article by default; set \CLS for another one.  With
%% pTeX the class option dvipdfmx is given.
%%
%% What is covered:
%%   * a box that starts in the middle of a page and is broken across
%%     several pages, with footnotes of one and of several lines in its
%%     title, in every part of its upper part and in its lower part;
%%   * a box with its only footnote in its last part;
%%   * a box with "height fixed for=first and middle": its parts fill
%%     the pages down to their footnotes;
%%   * a box with "page footnote" and an ordinary box inside a broken
%%     box.
%%
%% What to check:
%%   * TPF-RESULT OK (every footnote text is on the page of its mark),
%%     TPF-COUNTER OK, and no error or warning;
%%   * no "Overfull \vbox" and no "% split" (of \tracingpages: TeX had
%%     to split a footnote) in the log: the parts leave room for their
%%     footnotes;
%%   * more than six pages (the boxes are broken).
\ifdefined\CLS\else\def\CLS{article}\fi
\ifdefined\ptexversion\def\CLSOPT{dvipdfmx}\else\def\CLSOPT{}\fi
\expandafter\documentclass\expandafter[\CLSOPT]{\CLS}
\usepackage{tcbpagefootnote}
\tcbuselibrary{breakable}
\ifdefined\TPFhyperref\usepackage{hyperref}\fi
%% test-common.tex
%% ===============
%% Helpers shared by the test files; \input this file after
%% \usepackage{tcbpagefootnote}.
%%
%% \TPFfootnote[<kind>]{<text>}
%%   \footnote{<text>} with a check of where its text is typeset.
%%   <kind> (default: page) is one of
%%     page   a page footnote on the page of its mark;
%%     box    a footnote of the box (as in a minipage) on the page of
%%            its mark;
%%     later  a page footnote on a later page than its mark (a box
%%            broken across pages).
%% \TPFmark{<id>} and \TPFtext{<id>}{<kind>}
%%   The two halves of \TPFfootnote, for \footnote[<num>],
%%   \footnotemark and \footnotetext: put \TPFmark right after the mark
%%   and \TPFtext at the beginning of the text.  <id> is any name that
%%   is not a number (numbers are used by \TPFfootnote).
%% \TPFcounter{<counter>}{<value>}
%%   Check the value of <counter> at this point.
%% \TPFwidth (optional, defined before this file is input)
%%   A width, e.g. \def\TPFwidth{.75\columnwidth}: code is put around
%%   \@footnotetext of the page, as a package for a multi-column
%%   layout does, that sets the footnotes of the page in this width
%%   (computed from the \columnwidth that \@footnotetext is called
%%   with).  A page footnote is then a footnote set in this width.
%%
%% For each mark and each text, the page, the mark (\@thefnmark), the
%% hyperref anchor (if any) and, for a text, whether it is set in the
%% full \columnwidth of the page, or in \TPFwidth (a page footnote)
%% are written to the .aux file at shipout.  When the
%% .aux file is read back at \end{document}, every mark is checked:
%%   * its text is typeset exactly once, with the same mark;
%%   * the text is where <kind> says;
%%   * with hyperref, the anchor at the text is that of the mark.
%% Logs "TPF-NOTE NG: ..." for each failure, then "TPF-RESULT OK (<n>
%% notes)" or "TPF-RESULT NG".  A single run is enough.
\makeatletter
\ExplSyntaxOn
\int_new:N  \g__tpf_id_int
\seq_new:N  \g__tpf_ids_seq
\bool_new:N \g__tpf_ng_bool
\dim_new:N  \g__tpf_colwd_dim
\tl_new:N   \l__tpf_tmp_tl
\AddToHook { begindocument/end }
  {
    \cs_if_exist:NTF \TPFwidth
      {
        \dim_gset:Nn \g__tpf_colwd_dim { \TPFwidth }
        \cs_set_eq:NN \__tpf_footnotetext:n \@footnotetext
        \cs_set_protected:Npn \@footnotetext #1
          {
            \group_begin:
              \dim_set:Nn \columnwidth { \TPFwidth }
              \__tpf_footnotetext:n {#1}
            \group_end:
          }
      }
      { \dim_gset:Nn \g__tpf_colwd_dim { \columnwidth } }
  }

\cs_new_protected:Npn \__tpf_ng:n #1
  {
    \bool_gset_true:N \g__tpf_ng_bool
    \iow_term:e { TPF-NOTE~NG:~#1 }
  }

%% Write to the .aux file at shipout; #1 is frozen now except \thepage
\cs_new_protected:Npn \__tpf_shipout:n #1
  {
    \tl_set:Ne \l__tpf_tmp_tl {#1}
    \exp_args:NNV \iow_shipout_e:Nn \@auxout \l__tpf_tmp_tl
  }
\cs_new:Npn \__tpf_frozen_mark:
  { \exp_not:N \exp_not:n { \exp_not:o { \@thefnmark } } }
%% hyperref: the anchor of the mark, and the anchor at the text
%% (hyperref sets \@currentHref to the anchor it puts there).
%% Not with the new footnote code of LaTeX (latex-lab-footnotes), which
%% makes the links itself and does not use \Hy@footnote@currentHref.
\cs_new:Npn \__tpf_frozen_href:N #1
  {
    \bool_lazy_and:nnT
      { \cs_if_exist_p:N \Hy@footnote@currentHref }
      { ! \cs_if_exist_p:N \fnote_footnotetext:n }
      { \exp_not:N \exp_not:n { \exp_not:o {#1} } }
  }

\NewDocumentCommand \TPFmark { m }
  {
    \seq_gput_right:Nn \g__tpf_ids_seq {#1}
    \__tpf_shipout:n
      {
        \exp_not:N \exp_not:N \exp_not:N \TPF@m
        {#1} { \exp_not:N \thepage } { \__tpf_frozen_mark: }
        { \__tpf_frozen_href:N \Hy@footnote@currentHref }
      }
  }
\NewDocumentCommand \TPFtext { m m }
  {
    \__tpf_shipout:n
      {
        \exp_not:N \exp_not:N \exp_not:N \TPF@t
        {#1} {#2} { \exp_not:N \thepage } { \__tpf_frozen_mark: }
        { \dim_compare:nNnTF { \hsize } = { \g__tpf_colwd_dim } { 1 } { 0 } }
        { \__tpf_frozen_href:N \@currentHref }
      }
  }
\NewDocumentCommand \TPFfootnote { O{page} +m }
  {
    \int_gincr:N \g__tpf_id_int
    \exp_args:NV \__tpf_footnote:nnn \g__tpf_id_int {#1} {#2}
  }
\cs_new_protected:Npn \__tpf_footnote:nnn #1#2#3
  { \footnote { \TPFtext {#1} {#2} #3 } \TPFmark {#1} }

\NewDocumentCommand \TPFcounter { m m }
  {
    \int_compare:nNnTF { \value {#1} } = {#2}
      { \iow_term:e { TPF-COUNTER~OK:~#1=#2 } }
      { \__tpf_ng:n { counter~#1=\int_use:c { c@#1 },~want~#2 } }
  }

%% The records in the .aux file are ignored when it is read at
%% \begin{document} (they come from the previous run)
\cs_new_eq:NN \TPF@m \use_none:nnnn
\cs_new_eq:NN \TPF@t \use_none:nnnnnn
\cs_new_protected:Npn \__tpf_record_mark:nnnn #1#2#3#4
  {
    \tl_gset:cn { g__tpf_mhref_#1_tl } {#4}
    \tl_gset:cn { g__tpf_mpage_#1_tl } {#2}
    \tl_gset:cn { g__tpf_mmark_#1_tl } {#3}
  }
\cs_new_protected:Npn \__tpf_record_text:nnnnnn #1#2#3#4#5#6
  {
    \tl_gset:cn { g__tpf_thref_#1_tl } {#6}
    \int_if_exist:cF { g__tpf_tcount_#1_int }
      { \int_new:c { g__tpf_tcount_#1_int } }
    \int_gincr:c { g__tpf_tcount_#1_int }
    \tl_gset:cn { g__tpf_tkind_#1_tl } {#2}
    \tl_gset:cn { g__tpf_tpage_#1_tl } {#3}
    \tl_gset:cn { g__tpf_tmark_#1_tl } {#4}
    \tl_gset:cn { g__tpf_twide_#1_tl } {#5}
  }
\AddToHook { enddocument/afterlastpage }
  {
    \cs_gset_eq:NN \TPF@m \__tpf_record_mark:nnnn
    \cs_gset_eq:NN \TPF@t \__tpf_record_text:nnnnnn
  }
\AddToHook { enddocument/afteraux } { \__tpf_check: }

\cs_new_protected:Npn \__tpf_check:
  {
    \seq_map_function:NN \g__tpf_ids_seq \__tpf_check_one:n
    \seq_if_empty:NT \g__tpf_ids_seq { \__tpf_ng:n { no~footnotes } }
    \bool_if:NTF \g__tpf_ng_bool
      { \iow_term:n { TPF-RESULT~NG } }
      {
        \iow_term:e
          { TPF-RESULT~OK~(\seq_count:N \g__tpf_ids_seq\c_space_tl notes) }
      }
  }
\cs_new_protected:Npn \__tpf_check_one:n #1
  {
    \tl_if_exist:cTF { g__tpf_mpage_#1_tl }
      {
        \int_if_exist:cTF { g__tpf_tcount_#1_int }
          {
            \int_compare:nNnTF { \int_use:c { g__tpf_tcount_#1_int } } = { 1 }
              { \__tpf_check_text:n {#1} }
              {
                \__tpf_ng:n
                  { #1:~text~typeset~\int_use:c { g__tpf_tcount_#1_int }~times }
              }
          }
          { \__tpf_ng:n { #1:~text~lost } }
      }
      { \__tpf_ng:n { #1:~mark~not~shipped~out } }
  }
\cs_new_protected:Npn \__tpf_check_text:n #1
  {
    \tl_set:Ne \l__tpf_tmp_tl
      {
        #1:~\tl_use:c { g__tpf_tkind_#1_tl },~
        mark~\tl_use:c { g__tpf_mmark_#1_tl }~
        on~p.\tl_use:c { g__tpf_mpage_#1_tl },~
        text~\tl_use:c { g__tpf_tmark_#1_tl }~
        on~p.\tl_use:c { g__tpf_tpage_#1_tl }~
        \str_if_eq:eeTF { \tl_use:c { g__tpf_twide_#1_tl } } { 1 }
          { (page) } { (box) }
        \tl_if_empty:cF { g__tpf_mhref_#1_tl }
          {
            ,~anchor~\tl_use:c { g__tpf_mhref_#1_tl }~
            ->~\tl_use:c { g__tpf_thref_#1_tl }
          }
      }
    \bool_if:nTF
      {
        \tl_if_eq_p:cc { g__tpf_mmark_#1_tl } { g__tpf_tmark_#1_tl }
        && \__tpf_where_p:n {#1}
        && \tl_if_eq_p:cc { g__tpf_mhref_#1_tl } { g__tpf_thref_#1_tl }
      }
      { \iow_term:e { TPF-NOTE~OK:~\l__tpf_tmp_tl } }
      { \__tpf_ng:n { \l__tpf_tmp_tl } }
  }
\prg_new_conditional:Npnn \__tpf_where:n #1 { p }
  {
    \str_case_e:nnF { \tl_use:c { g__tpf_tkind_#1_tl } }
      {
        { page }
          {
            \bool_if:nTF
              {
                \int_compare_p:nNn { \tl_use:c { g__tpf_tpage_#1_tl } }
                  = { \tl_use:c { g__tpf_mpage_#1_tl } }
                && \str_if_eq_p:ee { \tl_use:c { g__tpf_twide_#1_tl } } { 1 }
              }
              { \prg_return_true: } { \prg_return_false: }
          }
        { box }
          {
            \bool_if:nTF
              {
                \int_compare_p:nNn { \tl_use:c { g__tpf_tpage_#1_tl } }
                  = { \tl_use:c { g__tpf_mpage_#1_tl } }
                && \str_if_eq_p:ee { \tl_use:c { g__tpf_twide_#1_tl } } { 0 }
              }
              { \prg_return_true: } { \prg_return_false: }
          }
        { later }
          {
            \bool_if:nTF
              {
                \int_compare_p:nNn { \tl_use:c { g__tpf_tpage_#1_tl } }
                  > { \tl_use:c { g__tpf_mpage_#1_tl } }
                && \str_if_eq_p:ee { \tl_use:c { g__tpf_twide_#1_tl } } { 1 }
              }
              { \prg_return_true: } { \prg_return_false: }
          }
      }
      { \prg_return_false: }
  }
\ExplSyntaxOff

%% \TPFfiller{<n>}: a paragraph of <n> filler sentences
\newcommand\TPF@filler{%
  This is a filler sentence of plain English text for the test. }
\newcommand\TPFfiller[1]{%
  \count@=0
  \loop\advance\count@ by 1 \TPF@filler\ifnum\count@<#1 \repeat\par}
\makeatother
\endinput

%% \TPFchunk{<n>}: <n> paragraphs, each followed by a line with a
%% footnote; every third footnote has several lines
\ExplSyntaxOn
\NewDocumentCommand \TPFchunk { m }
  {
    \int_step_inline:nn {#1}
      {
        \TPFfiller { 12 }
        \int_compare:nNnTF { \int_mod:nn {##1} { 3 } } = { 0 }
          { \TPFlongnote } { \TPFshortnote }
      }
  }
\ExplSyntaxOff
\newcommand\TPFshortnote{%
  A line with a footnote\TPFfootnote{A footnote of one line.}.\par}
\newcommand\TPFlongnote{%
  A line with a footnote\TPFfootnote{A footnote of several lines.
    This is a filler sentence of plain English text for the test.
    This is a filler sentence of plain English text for the test.
    This is a filler sentence of plain English text for the test.}.\par}
\begin{document}
\tracingpages=1
\section*{Breakable boxes broken by TeX}

\TPFfiller{20}
Body text\TPFfootnote{A footnote in the body text before the box.}.

\begin{tcolorbox}[breakable, page footnote,
  title={A long box\TPFfootnote{A footnote of the title.}}]
\TPFchunk{12}
\tcblower
\TPFchunk{6}
\end{tcolorbox}

Body text\TPFfootnote{A footnote in the body text after the box.}.

\begin{tcolorbox}[breakable, page footnote]
\TPFfiller{90}
The only footnote of this box\TPFfootnote{A footnote of the last part.}.
\end{tcolorbox}

\begin{tcolorbox}[breakable, page footnote, height fixed for=first and middle]
\TPFchunk{9}
\end{tcolorbox}

\begin{tcolorbox}[breakable, page footnote]
\TPFchunk{2}
\begin{tcolorbox}[page footnote]
A box inside the broken box\TPFfootnote{A footnote of the box inside.}.
\end{tcolorbox}
\TPFchunk{2}
\begin{tcolorbox}
An ordinary box inside the broken box\TPFfootnote[box]{A footnote of the
  ordinary box inside.}.
\end{tcolorbox}
\TPFchunk{2}
\end{tcolorbox}

\TPFcounter{footnote}{38}
\ifnum\value{page}>6 \else
  \typeout{TPF-PAGES NG: the boxes are not broken}%
\fi
\end{document}
